\begin{table}[htbp]
\centering
\caption{County-level border effect and entry-cap interaction, 100-mile border definition (CO, OR, MT, NM; quarterly)}
\label{tab:county_border_capped_main_100mi}
\footnotesize
\begin{tabular}{lccccc}
\toprule
& \multicolumn{2}{c}{Border} & \multicolumn{2}{c}{Border $\times$ Capped} & \\
Outcome & Coef. (\%) & $p$ (cl.) & Coef. (\%) & $p$ (cl.) & $n$ \\
\midrule
Revenue per capita & $+0.744$ ($+110.4\%$) & 0.001 & $-0.527$ ($-40.9\%$) & 0.055 & 3191 \\
Revenue per store & $+0.626$ ($+87.0\%$) & 0.060 & $-0.386$ ($-32.1\%$) & 0.222 & 1720 \\
Dispensary density & $+0.118$ ($+12.5\%$) & 0.677 & $-0.112$ ($-10.6\%$) & 0.597 & 1720 \\
\bottomrule
\end{tabular}
\begin{minipage}{0.98\textwidth}
\vspace{0.3em}
\footnotesize
\textit{Notes:} Sample: CO, OR, MT, NM (Washington dropped). Border = county centroid within 100 miles of the
nonlegal-state line, with $\geq$1,000 nonlegal-state residents within that reach (loosened from the 75-mile main
definition in Table~\ref{tab:county_border_capped_main}; 88 of 189 counties qualify as border here vs.\ 56 under
75 miles). Capped is the state-year entry-cap indicator (CO/MT always capped, NM always free, OR switches capped
in 2022). State and state $\times$ quarter fixed effects via a single concatenated state-quarter key. County-clustered
standard errors.
\end{minipage}
\end{table}
